\begin{minipage}{\linewidth}
\noindent\textbf{1.}  -- \textbf {(B).} Регулировка зеркал при перевозке пассажира. Пассажир должен:  \textbf {Сесть на сиденье рядом с вами.} \par \textit {Вес пассажира влияет на подвеску мотоцикла и вашу позу, поэтому зеркала необходимо настраивать, когда вы оба сидите на байке.} \\
\end{minipage}

\begin{minipage}{\linewidth}
\noindent\textbf{2.}  -- \textbf {(C).} Пассажир должен:  \textbf {Держаться за талию водителя.} \par \textit {Держась за талию водителя (или за бедра/пояс), пассажир будет синхронно двигаться и наклоняться вместе с водителем, обеспечивая устойчивость байка. Даже если вам от этого неловко.} \\
\end{minipage}

\begin{minipage}{\linewidth}
\noindent\textbf{3.}  -- \textbf {(A).} Когда мотоцикл останавливается на перекрёстке, пассажир должен:  \textbf {Держать ноги на подножках.} \par \textit {Если пасажир будет двигаться или опустит ноги на землю это может нарушить равновесие мотоцикла и создать риск получения травмы, поэтому во время остановки только водитель должен контролировать баланс и равновесие мотоцикла.} \\
\end{minipage}

\begin{minipage}{\linewidth}
\noindent\textbf{4.}  -- \textbf {(C).} При езде с пассажиром:  \textbf {Разгон и торможение займут больше времени.} \par \textit {Даже небольшой пассажир составляет значительную долю от общей массы мотоцикла — гораздо больше, чем два или три пассажира в автомобиле. Из-за этого мотоциклом сложнее управлять, а разгон и торможение занимают больше времени. Езды с пасажиром лучше избегать в принципе.} \\
\end{minipage}

\begin{minipage}{\linewidth}
\noindent\textbf{5.}  -- \textbf {(A).} При езде с пассажирами вам следует:  \textbf {Начинать тормозить раньше чем обычно.} \par \textit {Даже небольшой пассажир составляет значительную долю от общей массы мотоцикла — гораздо большую, чем два или три пассажира в автомобиле. Это затрудняет управление мотоциклом и замедляет его разгон и торможение.} \\
\end{minipage}

\begin{minipage}{\linewidth}
\noindent\textbf{6.}  -- \textbf {(D).} При перевозке грузов:  \textbf {Располагайте груз низко и как можно ближе к передней части.} \par \textit {Низко расположенный и смещенный вперед груз позволяет сохранить исходный центр тяжести байка, обеспечивая стабильность и управляемость. Груз, расположенный позади задней оси, может вызвать вудди или занос.} \\
\end{minipage}

\begin{minipage}{\linewidth}
\noindent\textbf{7.}  -- \textbf {(D).} Группа байкеров переходит из шахматного порядка в колонну по одному, в случае:  \textbf {При прохождении поворотов.} \par \textit {Повороты, изгибы, а также съезд-въезд на автомагистраль требуют, чтобы у мотоциклиста было достаточно места как впереди так и по бокам. Для маневра и реагирования. Поэтому группа едет гуськом. А затем вновь пересобираивается в шахматную формацию.} \\
\end{minipage}

\begin{minipage}{\linewidth}
\noindent\textbf{8.}  -- \textbf {(D).} На этом изображении мотоциклисты едут в шахматном порядке. Дистанция между A и C должна составлять не менее:  \textbf {Две секунды.} \par \textit {Езда в шахматном порядке позволяет группе оставаться компактной, при этом каждый мотоциклист сохраняет 2-х секундный интервал до мотоцикла непосредственно впереди, что дает пространство для реакции и маневра. Это минимум. Реальная дистанция должна соответствовать вашим навыкам и обстановке.} \\
\end{minipage}

\begin{minipage}{\linewidth}
\noindent\textbf{9.}  -- \textbf {(C).} Мотоциклисты, едущие в шахматном порядке, будут обгонять автомобиль. После того как лидер закончит обгон, он должен:  \textbf {Занять левую часть полосы движения.} \par \textit {Мотоциклисты обгоняют по одному, даже в при езде в шахматной колонне. После обгона ведущий мотоциклист встает в левую часть полосы движения где он и был до обгона, проезжает достаточно вперед освобождая место для следующего, который выходит вперед в качестве нового лидера и обгоняет автомобиль таким же образом. Повторить пока вся группа не проедет и не восстановит прежний шахматный порядок перед обгоняемым авто.} \\
\end{minipage}

\begin{minipage}{\linewidth}
\noindent\textbf{10.}  -- \textbf {(A).} Алкоголь попадает в кровоток и быстро достигает мозга. На что это влияет в первую очередь?  \textbf {Способность принимать адекватные решения.} \par \textit {Алкоголь в первую очередь ухудшает способность к правильной оценке ситуации и ослабляет сдерживающие факторы, поэтому водитель может идти на риск и неправильно оценивать степень своего опьянения до того, как проявятся физические симптомы, такие как невнятная речь, нарушение равновесия или проблемы со зрением.} \\
\end{minipage}

\begin{minipage}{\linewidth}
\noindent\textbf{11.}  -- \textbf {(C).} Алкоголь в первую очередь влияет на:  \textbf {Способность принимать адекватные решения.} \par \textit {Алкоголь в первую очередь ухудшает способность к правильной оценке ситуации и ослабляет сдерживающие факторы, поэтому мотоциклист может идти на риск с большей легкостью и неправильно оценивать степень своего опьянения до того, как проявятся такие физические симптомы, как невнятная речь, нарушение равновесия или проблемы со зрением.} \\
\end{minipage}

\begin{minipage}{\linewidth}
\noindent\textbf{12.}  -- \textbf {(C).} Примерно половина всех смертельных случаев в результате ДТП с участием мотоциклов происходит по причине:  \textbf {Из-за воздействия алкоголя.} \par \textit {Алкоголь ухудшает способность к правильной оценке ситуации и ослабляет сдерживающие факторы, поэтому мотоциклист может идти на риск и неправильно оценивать степень своего опьянения до появления таких физических симптомов, как невнятная речь, нарушение равновесия или проблемы со зрением.} \\
\end{minipage}

\begin{minipage}{\linewidth}
\noindent\textbf{13.}  -- \textbf {(B).} Если вы не контролируете количество выпитого:  \textbf {Перед поездкой выдержите как минимум час после каждой порции алкоголя.} \par \textit {Только время выводит алкоголь из организма — примерно по одной порции алкоголя в час (что означает порция см. раздел про алкоголь, описание термина BAC). Езда на более низкой скорости, физические упражнения или кофе не снижают степень опьянения.} \\
\end{minipage}

\begin{minipage}{\linewidth}
\noindent\textbf{14.}  -- \textbf {(B).} Исследования показывают, что почти половина всех мотоциклистов, погибших в ДТП:  \textbf {Употребляли алкоголь.} \par \textit {Алкоголь является фактором в значительной части смертельных ДТП с участием мотоциклов. Не садитесь за руль в нетрезвом состоянии.} \\
\end{minipage}

\begin{minipage}{\linewidth}
\noindent\textbf{15.}  -- \textbf {(A).} Какое количество алкоголя начинает влять на способность ясно мыслить и безопасно управлять мотоциклом?  \textbf {1 порция алкоголя.} \par \textit {Даже одна порция алкоголя может замедлить реакцию и ухудшить способность принимать ответственные, адекватные, взвешенные решения, поэтому для мотоциклиста не существует безопасного количества алкоголя. Не садитесь за руль в нетрезвом состоянии.} \\
\end{minipage}

\begin{minipage}{\linewidth}
\noindent\textbf{16.}  -- \textbf {(A).} Если вы выпьете 4 стандартных порции алкоголя, какое минимальное количество алкоголя останется в вашем организме через 2 часа?  \textbf {2 порции.} \par \textit {Организм выводит примерно одну порцию алкоголя в час, поэтому через 2 часа из 4 порций исчезнут не более 2, а в организме останется не менее 2 (см. раздел про алкоголь, термин BAC).} \\
\end{minipage}

\begin{minipage}{\linewidth}
\noindent\textbf{17.}  -- \textbf {(C).} Для максимальной безопасности ваш шлем должен:  \textbf {Плотно прилегать по всему периметру.} \par \textit {Плотное и равномерное прилегание предотвращает смещение или срыв шлема при аварии. Шлемы спасают жизни. Носите полнолицевой шлем. Не искушайте судьбу. Не рискуйте.} \\
\end{minipage}

\begin{minipage}{\linewidth}
\noindent\textbf{18.}  -- \textbf {(B).} Самое важное, что вы можете сделать, чтобы повысить свои шансы на выживание при аварии это:  \textbf {Носить надежно закрепленный шлем.} \par \textit {Травмы головы являются основной причиной смертности среди мотоциклистов, а правильно закрепленный шлем — самая эффективная защита от них. Шлемы спасают жизни. Носите полнолицевой шлем. Не искушайте судьбу. Не рискуйте.} \\
\end{minipage}

\begin{minipage}{\linewidth}
\noindent\textbf{19.}  -- \textbf {(C).} Ношение шлема может снизить вероятность получения смертельной травмы головы при:  \textbf {Любом ДТП, независимо от скорости.} \par \textit {Шлем снижает риск смертельной травмы головы при любом типе ДТП. Шлемы спасают жизни. Носите полнолицевой шлем. Не искушайте судьбу. Не рискуйте.} \\
\end{minipage}

\begin{minipage}{\linewidth}
\noindent\textbf{20.}  -- \textbf {(A).} Ношение шлема:  \textbf {Снизит вероятность травм шеи и головы.} \par \textit {Шлем не вызывает травм шеи и не ограничивает обзор. Все это — городские легенды. ОБС. Носите полнолицевой шлем. Не искушайте судьбу. Не рискуйте.} \\
\end{minipage}

\begin{minipage}{\linewidth}
\noindent\textbf{21.}  -- \textbf {(C).} Что из перечисленного ниже обеспечивает наилучшую защиту глаз и лица во время езды?  \textbf {Защитный щиток шлема.} \par \textit {Защитный щиток шлема полностью закрывает лицо и глаза, защищая их от ветра, мусора и ударов. Все остальные варианты обеспечивают лишь частичную защиту.} \\
\end{minipage}

\begin{minipage}{\linewidth}
\noindent\textbf{22.}  -- \textbf {(D).} При езде на незнакомом мотоцикле:  \textbf {Оставляйте больше дистанции для торможения.} \par \textit {Незнакомые тормоза и управляемость сложно оценить, поэтому оставляйте большее запас расстояния для остановки, пока не освоите мотоцикл.} \\
\end{minipage}

\begin{minipage}{\linewidth}
\noindent\textbf{23.}  -- \textbf {(C).} Аварийный выключатель двигателя обычно расположен:  \textbf {Рядом с правой ручкой} \par \textit {Аварийный выключатель двигателя — важный элемент системы безопасности. Представьте себе мотоцикл, скользящий на боку, при этом заднее колесо продолжает бешенно вращаться. Обычно кнопка находится рядом с правой ручкой, чтобы вы могли ее нажать одним движением большого пальца. Тренируйте рефлекс выключения — глушите мото на порковке с помощью Cut-off switch.} \\
\end{minipage}

\begin{minipage}{\linewidth}
\noindent\textbf{24.}  -- \textbf {(C).} При правильной регулировке зеркала вы можете видеть:  \textbf {Часть полосы движения рядом с вами и позади вас.} \par \textit {Зеркала мотоцикла маленькие и показывают только часть полосы движения сзади и сбоку от вас. Всегда остаются слепые зоны, и вам нужно поворачивать голову, чтобы проверить обстановку (Head-Checks).} \\
\end{minipage}

\begin{minipage}{\linewidth}
\noindent\textbf{25.}  -- \textbf {(C).} При замедлении перед обычной остановкой:  \textbf {При замедлении переключайтесь на более низкую передачу.} \par \textit {Переключение на более низкую передачу при замедлении позволяет оставаться на правильной передаче, если придется разгоняться, а к моменту полной остановки (например, на красном светофоре) вы уже будете на 1-й передаче, готовые  продолжить движение. Использование только одного тормоза или движение накатом с выжатым сцеплением снижает управляемость.} \\
\end{minipage}

\begin{minipage}{\linewidth}
\noindent\textbf{26.}  -- \textbf {(D).} При повороте следует:  \textbf {Смотреть вперед, в сторону поворота.} \par \textit {Мотоцикл имеет тенденцию двигаться туда, куда вы смотрите, поэтому смотрите вперед, в сторону поворота, туда, куда вы хотите ехать. Если вы будете фиксировать взгляд на центральной линии или ограждении, это потянет вас в их сторону. Смотреть надо на точку выхода из поворота, а следить за обстановкой боковым зрением.} \\
\end{minipage}

\begin{minipage}{\linewidth}
\noindent\textbf{27.}  -- \textbf {(B).} Какое из следующих утверждений верно:  \textbf {Нажми с левой стороны, наклонись влево, поедешь влево.} \par \textit {Технически, просто чтобы запомнить ответ: нажатие, наклон и направление движения всегда совпадают. А чтобы сделать поворот правильно надо потянуть подбородок к зеркалу. Это еще называют mirror kiss. Это автоматически сделает за вас три действия: правильный наклон тела вперед и в сторону, нажатие на правильныю рукоятнку с нужной силой и взгляд повер дороги. Одно простое действие. Легкий наклон вперед снимает страх свалиться с байка, а также вызывает рефлекс сжатия бака коленями. Это магия.} \\
\end{minipage}

\begin{minipage}{\linewidth}
\noindent\textbf{28.}  -- \textbf {(A).} Во время поворота лучше всего:  \textbf {Сохранять постоянную скорость или плавно ускоряться.} \par \textit {Равномерное или плавное увеличение газа позволяет сохранить стабильный угол наклона мотоцикла и сцепление шин с дорогой. Резкие ускорение, замедление или торможение в процессе поворота нарушают сцепление с дорогой, изменяют угол наклона мотоцикла и, следовательно, траекторию поворота.} \\
\end{minipage}

\begin{minipage}{\linewidth}
\noindent\textbf{29.}  -- \textbf {(C).} Если вам нужно замедлиться во время поворота:  \textbf {Аккуратно задействуйте оба тормоза.} \par \textit {Основное правило — всегда использовать оба тормоза. Резкое торможение может привести к заносу, но само по себе торможение в повороте не запрещено. Изменение в скорости или тяге в повороте меняет радиус.} \\
\end{minipage}

\begin{minipage}{\linewidth}
\noindent\textbf{30.}  -- \textbf {(B).} Для обычной остановки используйте:  \textbf {Оба тормоза и переключение на пониженную передачу.} \par \textit {Одновременное использование обоих тормозов обеспечивает максимальную тормозную силу и контроль, а переключение на пониженную передачу позволяет остановиться на 1-й передаче, готовым к возобновлению движения. Использование только одного тормоза снижает контроль.} \\
\end{minipage}

\begin{minipage}{\linewidth}
\noindent\textbf{31.}  -- \textbf {(C).} У вашего мотоцикла два тормоза. Используйте оба тормоза:  \textbf {Каждый раз, когда вы замедляетесь или останавливаетесь.} \par \textit {Передний тормоз обеспечивает большую часть тормозного усилия, но использование обоих тормозов каждый раз формирует привычку, которая понадобится вам в экстренной ситуации. Используйте оба тормоза каждый раз, когда вы замедляетесь или останавливаетесь.} \\
\end{minipage}

\begin{minipage}{\linewidth}
\noindent\textbf{32.}  -- \textbf {(B).} При обычном повороте:  \textbf {Наклоняйте мотоцикл и тело под одинаковым углом.} \par \textit {При нормальной скорости поворота мотоцикл и мотоциклист наклоняются вместе под одним и тем же углом, что позволяет сохранить равновесие мотоцикла. Это базовое правило. На курсах BRC это будет разобрано более подробно.} \\
\end{minipage}

\begin{minipage}{\linewidth}
\noindent\textbf{33.}  -- \textbf {(A).} При медленных крутых поворотах:  \textbf {Наклоняйте только мотоцикл, а тело держите прямо.} \par \textit {На низкой скорости сам байк должен наклоняться, чтобы совершать крутые повороты, но ваше тело, остающееся в вертикальном положении, действуя как противовес.} \\
\end{minipage}

\begin{minipage}{\linewidth}
\noindent\textbf{34.}  -- \textbf {(D).} Чтобы снизить скорость перед поворотом:  \textbf {Закройте дроссель и задействуйте оба тормоза.} \par \textit {Снизьте скорость перед поворотом. Никогда не знаешь, насколько хорошее сцепление с дорогой. Выработайте привычку тормозить так: сбросить газ, задействовать оба тормоза.} \\
\end{minipage}

\begin{minipage}{\linewidth}
\noindent\textbf{35.}  -- \textbf {(D).} При повороте важно:  \textbf {Смотреть сквозь поворот туда, куда вы хотите ехать.} \par \textit {Мотоцикл имеет тенденцию двигаться туда, куда вы смотрите, поэтому смотрите сквозь поворот на то место, куда вы хотите выехать. Замедление следует производить до поворота, а не во время него или при выходе из него. Вы никогда не знаете, насколько хорошее сцепление с дорогой у вас будет в повороте.} \\
\end{minipage}

\begin{minipage}{\linewidth}
\noindent\textbf{36.}  -- \textbf {(B).} Важно искать на дороге пути отхода:  \textbf {На перекрестках.} \par \textit {Перекрестки, вероятно, являются самыми опасными местами, поэтому всегда прикидывайте, куда можно свернуть или уклониться, если что-то пойдет не так.} \\
\end{minipage}

\begin{minipage}{\linewidth}
\noindent\textbf{37.}  -- \textbf {(B).} Чтобы правильно принимать решения в дорожном движении, сначала нужно:  \textbf {Смотреть вперёд.} \par \textit {Правильное принятие решений начинается с информации: сначала осмотрите дорогу впереди, затем оцените то, что видите, и потом действуйте соответственно. Методика SEE (Search, Evaluate,Execute).} \\
\end{minipage}

\begin{minipage}{\linewidth}
\noindent\textbf{38.}  -- \textbf {(B).} Водитель автомобиля подаёт сигнал о повороте налево. Мотоциклист должен:  \textbf {Снизить скорость и держаться подальше от приближающихся транспортных средств.} \par \textit {Даже если у вас право проехать первым, водитель автомобиля может неправильно оценить вашу скорость или вообще не заметить вас. Снизте скорость и держитесь он него на безопасной дистанции.} \\
\end{minipage}

\begin{minipage}{\linewidth}
\noindent\textbf{39.}  -- \textbf {(D).} Вам следует уделять особое внимание проверке обстановки сзади:  \textbf {Перед тем, как снизить скорость.} \par \textit {Силуэт байка узкий, малозаметный. Оглянитесь назад перед тем, как снизить скорость, чтобы убедиться, что следующий за вами водитель видит вас и сможет вовремя остановиться. При необходимости вы можете дополнительно предупредить его, мигая стоп-сигналом.} \\
\end{minipage}

\begin{minipage}{\linewidth}
\noindent\textbf{40.}  -- \textbf {(A).} При подготовке к обгону обычно лучше:  \textbf {Подать сигнал, посмотреть в левое зеркало,  посмотреть налево повернув голову.} \par \textit {Каверзный вопрос о повороте головы (Head-checks). Сначала включается поворотник — это очевидно, затем смотрим в зеркало, потому что это просто, а потом поворачиваем голову и смотрим глазами, что ситуация не изменилась и никого нет в слепых зонах. То что доктор прописал.} \\
\end{minipage}

\begin{minipage}{\linewidth}
\noindent\textbf{41.}  -- \textbf {(C).} Движение по центру полосы:  \textbf {Не позволяет другим транспортным средствам занимать вашу полосу.} \par \textit {Центральная часть полосы делает вас более заметным и не даёт другому транспортному средству вклиниться в вашу полосу и ехать рядом с вами. О движении по центру полосы существует множество мифов. Внимание: хотя это лучшее решение, у него есть свои недостатки!} \\
\end{minipage}

\begin{minipage}{\linewidth}
\noindent\textbf{42.}  -- \textbf {(B).} Чтобы предотвратить попытки других водителей занять вашу полосу вместе с вами, обычно лучше всего:  \textbf {Двигаться по центральной части полосы.} \par \textit {Движение по центральной части полосы визуально занимает всю полосу, поэтому у других водителей меньше искушения обогнать вас или занять полосу рядом. Смещение к любому краю провоцирует их на это. О движении по центру полосы существует множество мифов. Обратите внимание: хотя это лучшее решение, у него есть свои недостатки.} \\
\end{minipage}

\begin{minipage}{\linewidth}
\noindent\textbf{43.}  -- \textbf {(B).} Поддержание зоны безопасности (пространственного буфера) вокруг мотоцикла:  \textbf {Даёт вам время и пространство для реакции.} \par \textit {Свободное пространство вокруг даёт время заметить опасность и место для торможения или манёвра. Это необходимо каждому мотоциклисту.} \\
\end{minipage}

\begin{minipage}{\linewidth}
\noindent\textbf{44.}  -- \textbf {(A).} Чтобы вовремя заметить встречный транспорт и сохранить зону безопасности, как показано на иллюстрации выше:  \textbf {Двигайтесь по центру полосы.} \par \textit {При прохождении извилистого поворота опасно приближаться к разделительной линии, так как встречный транспорт может срезать угол и выехать на вашу полосу. Правая сторона грязная и обеспечивает плохой обзор впереди. Центр полосы оставляет запас места с обеих сторон, сохраняя буфер как от встречного транспорта, так и от обочины.} \\
\end{minipage}

\begin{minipage}{\linewidth}
\noindent\textbf{45.}  -- \textbf {(B).} Вы планируете повернуть направо после остановки на перекрёстке. Обычно лучше всего:  \textbf {Выполнять поворот из центральной части полосы.} \par \textit {Нахождение в центральной части не позволяет другим транспортным средствам втиснуться мимо вас и проехать по вашей же полосе во время поворота. О движении по центру полосы существует множество мифов. Обратите внимание: хотя это лучшее решение, у него есть свои недостатки.} \\
\end{minipage}

\begin{minipage}{\linewidth}
\noindent\textbf{46.}  -- \textbf {(D).} Движение на близкой дистанции за впереди идущим транспортом:  \textbf {Является одной из причин ДТП по вине мотоциклистов.} \par \textit {Несоблюдение дистанции не оставляет времени и места для реакции или остановки, а также перекрывает обзор дороги и дорожного покрытия.} \\
\end{minipage}

\begin{minipage}{\linewidth}
\noindent\textbf{47.}  -- \textbf {(A).} Двигаясь за автомобилем, вам следует стараться:  \textbf {Ехать так, чтобы водитель видел вас в зеркало заднего вида.} \par \textit {Водители обычно смотрят только в зеркало заднего вида. Поэтому центральная часть полосы — ваш лучший выбор.} \\
\end{minipage}

\begin{minipage}{\linewidth}
\noindent\textbf{48.}  -- \textbf {(A).} Самой частой причиной ДТП с участием мотоциклов и автомобилей на перекрёстках является:  \textbf {Нарушение водителями преимущественного права проезда мотоциклиста.} \par \textit {Самый распространённый тип аварии — когда водитель не замечает мотоцикл и выезжает или поворачивает прямо перед мотоциклистом.} \\
\end{minipage}

\begin{minipage}{\linewidth}
\noindent\textbf{49.}  -- \textbf {(A).} ДТП с участием автомобилей и мотоциклов чаще всего происходят:  \textbf {На перекрёстках.} \par \textit {Перекрёстки — это места, где пересекаются траектории движения, а водители чаще всего просто не замечают мотоциклистов.} \\
\end{minipage}

\begin{minipage}{\linewidth}
\noindent\textbf{50.}  -- \textbf {(C).} Зрительный контакт с водителем на перекрёстке:  \textbf {Не означает, что водитель действительно вас видит.} \par \textit {Водитель может смотреть прямо на вас, но при этом не осознавать присутствие мотоцикла. Поэтому зрительный контакт не гарантирует, что вас заметили и уступят дорогу. Будьте готовы реагировать.} \\
\end{minipage}

\begin{minipage}{\linewidth}
\noindent\textbf{51.}  -- \textbf {(B).} Наибольшая опасность для мотоциклиста при проезде мимо припаркованных автомобилей — это:  \textbf {Автомобиль, выезжающий с парковки на дорогу.} \par \textit {Водитель, отъезжающий от бордюра, может перекрыть всю улицу, не оставив вам путей для отступления. Это наибольшая опасность, поэтому будьте внимательны и следите за водителями в припаркованных машинах.} \\
\end{minipage}

\begin{minipage}{\linewidth}
\noindent\textbf{52.}  -- \textbf {(D).} На этом рисунке световой сигнал, который скорее всего сделает мотоцикл заметным для водителя автомобиля — это:  \textbf {Указатель поворота.} \par \textit {Единственный полезный сигнал в этой ситуации — указатель поворота. Но имейте в виду, что водители всё равно его не замечают.} \\
\end{minipage}

\begin{minipage}{\linewidth}
\noindent\textbf{53.}  -- \textbf {(D).} При движении в потоке ночью лучший способ обнаружить неровности на дороге — это:  \textbf {Следить за задними габаритными огнями впереди идущего автомобиля.} \par \textit {Колебания задних огней вверх и вниз указывают на кочки или неровное покрытие впереди. Дальний свет мало помогает в плотном потоке и может ослепить других водителей.} \\
\end{minipage}

\begin{minipage}{\linewidth}
\noindent\textbf{54.}  -- \textbf {(D).} Лучший способ сделать ваш мотоцикл заметнее для других — это:  \textbf {Ездить с включённой фарой.} \par \textit {Включённая фара делает вас более заметным для других водителей в любое время суток. Положение на полосе и цвет шлема помогают значительно меньше.} \\
\end{minipage}

\begin{minipage}{\linewidth}
\noindent\textbf{55.}  -- \textbf {(A).} Забытый включённый указатель поворота скорее всего:  \textbf {Спровоцирует других водителей выехать на вашу траекторию движения.} \par \textit {На мотоцикле поворотник не выключается автоматически. Включённый сигнал поворота ошибочно информирует окружающих о том, что вы собираетесь повернуть, поэтому водители могут выехать или повернуть прямо перед вами.} \\
\end{minipage}

\begin{minipage}{\linewidth}
\noindent\textbf{56.}  -- \textbf {(D).} Особенно важно проверять зеркала заднего вида:  \textbf {Перед замедлением или остановкой.} \par \textit {Мотоцикл узкий, и водители сзади могут не заметить, что он замедляется. Сначала проверьте зеркала, чтобы оценить дистанцию до едущего сзади транспорта, а затем поморгате стоп-сигналом, чтобы привлечь внимание перед замедлением или остановкой.} \\
\end{minipage}

\begin{minipage}{\linewidth}
\noindent\textbf{57.}  -- \textbf {(A).} Перед перестроением на левую полосу:  \textbf {Проверьте левое зеркало и поверните голову влево для проверки слепой зоны.} \par \textit {У зеркал есть слепые зоны. Поворот головы — обязательное правило: поверните голову и посмотрите в нужную сторону перед перестроением.} \\
\end{minipage}

\begin{minipage}{\linewidth}
\noindent\textbf{58.}  -- \textbf {(C).} Лучшая экипировка, чтобы быть заметным ночью — это:  \textbf {Одежда со световозвращающими элементами.} \par \textit {Световозвращающий материал отражает свет фар. Это лучший выбор для ночной езды.} \\
\end{minipage}

\begin{minipage}{\linewidth}
\noindent\textbf{59.}  -- \textbf {(B).} Чтобы лучше видеть дорогу при движении за автомобилем ночью:  \textbf {Используйте свет фар впереди идущего автомобиля.} \par \textit {Фары едущего впереди автомобиля мощнее ваших и освещают дорогу для вас, позволяя видеть дальше и вовремя замечать опасности.} \\
\end{minipage}

\begin{minipage}{\linewidth}
\noindent\textbf{60.}  -- \textbf {(C).} При езде в ночное время:  \textbf {Носите световозвращающую одежду.} \par \textit {Световозвращающая одежда делает вас намного более заметным в темноте.} \\
\end{minipage}

\begin{minipage}{\linewidth}
\noindent\textbf{61.}  -- \textbf {(D).} Если при остановке вы заблокировали переднее колесо:  \textbf {Отпустить и снова нажать передний тормоз.} \par \textit {Вопрорс с подвохом. Для переднего и заднего тормозов действуют разные базовые правила. Заблокированное переднее колесо очень опасно: вы теряете управление и баланс, и мотоцикл может упасть. Немедленно отпустите передний тормоз, чтобы восстановить сцепление с дорогой, а затем снова нажмите его с меньшим усилием.} \\
\end{minipage}

\begin{minipage}{\linewidth}
\noindent\textbf{62.}  -- \textbf {(D).} На рисунке выше при попытке экстренной остановки заблокировалось заднее колесо. Как правило, лучше всего:  \textbf {Держать заднее колесо заблокированным до полной остановки.} \par \textit {Для переднего и заднего тормозов действуют разные правила. Заблокированное заднее колесо должно оставаться заблокированным до полной остановки. Высок риск поймать хайсайд (highside). Это базовое правило. Оно не идеальное. На BRC учат как с этим справиться.} \\
\end{minipage}

\begin{minipage}{\linewidth}
\noindent\textbf{63.}  -- \textbf {(B).} При резком объезде препятствия важно:  \textbf {Разделять торможение и маневр.} \par \textit {Торможение и маневр требуют сцепления шин с дорогой, поэтому их одновременное выполнение может привести к заносу. Сначала объезжайте, тормозите, когда мотоцикл выровняется, либо сначала затормозите, а затем отпустите тормоз и делайте маневр уклонения.} \\
\end{minipage}

\begin{minipage}{\linewidth}
\noindent\textbf{64.}  -- \textbf {(A).} Для проезда по металлическим решетчатым настилам моста:  \textbf {Двигаться прямо.} \par \textit {Металлические решетки вызывают ощущение неустойчивости мотоцикла, но шины сохраняют сцепление. Расслабьтесь, удерживайте постоянную скорость и двигайтесь прямо. Движение зигзагом или под углом только усугубит ситуацию.} \\
\end{minipage}

\begin{minipage}{\linewidth}
\noindent\textbf{65.}  -- \textbf {(B).} При переезде через предмет на дороге:  \textbf {Приподняться над сиденьем.} \par \textit {Небольшой подъем над сиденьем переносит ваш вес на подножки, что понижает центр тяжести и улучшает баланс. Кроме того, согнутые в коленях ноги амортизируют удар вместо позвоночника.} \\
\end{minipage}

\begin{minipage}{\linewidth}
\noindent\textbf{66.}  -- \textbf {(B).} При проезде препятствия лучше всего:  \textbf {Приподняться над сиденьем.} \par \textit {Положение стоя на подножках со слегка согнутыми коленями позволяет мотоциклу двигаться под вами. Это сохраняет низкий центр тяжести, гасит удар от препятствия и помогает сохранить контроль.} \\
\end{minipage}

\begin{minipage}{\linewidth}
\noindent\textbf{67.}  -- \textbf {(C).} На двухполосной дороге с твердым покрытием песок и гравий чаще всего скапливаются:  \textbf {У краев дороги (у обочин).} \par \textit {Поток автомобилей выталкивает песок и гравий из автомобильной колеи к краям дороги, поэтому они скапливаются ближе к обочине.} \\
\end{minipage}

\begin{minipage}{\linewidth}
\noindent\textbf{68.}  -- \textbf {(D).} На мокрой дороге наиболее скользко обычно:  \textbf {В центральной части полосы.} \par \textit {Обычно центр полосы — лучшее место для движения, но не во время дождя. Масло и технические жидкости капают именно туда и смешиваются с водой, поэтому центр становится очень скользким. Двигайтесь по автомобильной колее.} \\
\end{minipage}

\begin{minipage}{\linewidth}
\noindent\textbf{69.}  -- \textbf {(B).} Чтобы безопасно проехать скользкий участок дороги, лучше всего:  \textbf {Снизить скорость до въезда на него.} \par \textit {Торможение, ускорение и переключение передач вниз создают дополнительную нагрузку на сцепление шин. Снизьте скорость до скользкого участка, затем удерживайте мотоцикл ровно и стабильно при его проезде.} \\
\end{minipage}

\begin{minipage}{\linewidth}
\noindent\textbf{70.}  -- \textbf {(D).} Для остановки на мокром асфальте:  \textbf {Использовать передний и задний тормоза одновременно.} \par \textit {Использование обоих тормозов обеспечивает наилучшее торможение даже на мокрой дороге. Нажимайте на них плавно и мягко, так как на мокром асфальте колесо легче заблокировать.} \\
\end{minipage}

\begin{minipage}{\linewidth}
\noindent\textbf{71.}  -- \textbf {(D).} При движении по дорогам с твердым покрытием в местах скопления песка и гравия следует:  \textbf {Избегать резких изменений скорости или направления движения.} \par \textit {Сыпучие материалы обеспечивают слабое сцепление. Резкое торможение, ускорение или маневрирование могут вызвать занос, поэтому рулите плавно и степенно.} \\
\end{minipage}

\begin{minipage}{\linewidth}
\noindent\textbf{72.}  -- \textbf {(C).} Если спущено заднее колесо:  \textbf {Задняя часть мотоцикла будет вилять из стороны в сторону.} \par \textit {Спущенное заднее колесо заставляет заднюю часть мотоцикла раскачиваться или вилять из стороны в сторону. Вы все еще можете сохранить контроль, крепко удерживая руль.} \\
\end{minipage}

\begin{minipage}{\linewidth}
\noindent\textbf{73.}  -- \textbf {(B).} Если во время движения у вас пробило колесо, крепко держитесь за руль и:  \textbf {Используйте тормоз исправного колеса или избегайте торможения.} \par \textit {Торможение спущенным колесом может вызвать занос и падение. Если необходимо затормозить, аккуратно используйте тормоз исправного колеса или избегайте торможения вовсе, позволяя мотоциклу замедлиться самому.} \\
\end{minipage}

\begin{minipage}{\linewidth}
\noindent\textbf{74.}  -- \textbf {(C).} Если во время движения спущено переднее колесо:  \textbf {Используйте задний тормоз и съедьте к краю дороги.} \par \textit {Заднее колесо исправно, поэтому аккуратно задействуйте его тормоз. Никогда не тормозите спущенным колесом. Снизьте скорость и только затем съезжайте к краю дороги.} \\
\end{minipage}

\begin{minipage}{\linewidth}
\noindent\textbf{75.}  -- \textbf {(A).} При спущенном переднем колесе:  \textbf {Рулевое управление станет тяжелым.} \par \textit {Спущенное переднее колесо делает рулевое управление тяжелым и инертным. Виляние задней части из стороны в сторону — признак спущенного заднего колеса.} \\
\end{minipage}

\begin{minipage}{\linewidth}
\noindent\textbf{76.}  -- \textbf {(A).} Если заклинило ручку газа:  \textbf {Выжмите сцепление и используйте аварийный выключатель двигателя.} \par \textit {Выжатое сцепление не даст мотоциклу продолжать разгон, но заклинившая ручка газа раскрутит мотор до критических оборотов, что может его повредить. Аварийный выключатель двигателя заглушит мотор, позволив вам безопасно затормозить и съехать на обочину.} \\
\end{minipage}

\begin{minipage}{\linewidth}
\noindent\textbf{77.}  -- \textbf {(D).} На рисунке выше обычно лучше всего:  \textbf {Двигаться в позиции 2 или 3.} \par \textit {Позиции 2 и 3 позволяют держаться подальше от опасности в левой части полосы и оставляет запас для реакции.} \\
\end{minipage}

\begin{minipage}{\linewidth}
\noindent\textbf{78.}  -- \textbf {(D).} На этом рисунке A, B и C едут в группе. Для обгона мотоциклист B должен:  \textbf {Вернуться в правую часть полосы позади A.} \par \textit {Мотоциклисты в группе совершают обгон по одному. Каждый возвращается на свою исходную позицию в строю. Цель состоит в том, чтобы восстановить порядок колонны после обгона.} \\
\end{minipage}

\begin{minipage}{\linewidth}
\noindent\textbf{79.}  -- \textbf {(B).} Чтобы создать больше пространства в ситуации, изображенной выше, вам следует:  \textbf {Скорректировать скорость и оставаться по центру полосы.} \par \textit {Вливающийся в поток водитель, скорее всего, вас не заметит, поэтому дайте ему проехать. Да, вам обязаны уступить, но ваша безопасность важнее.} \\
\end{minipage}

\begin{minipage}{\linewidth}
\noindent\textbf{80.}  -- \textbf {(A).} Большинство аварий с участием мотоцикла и автомобиля:  \textbf {Происходят из-за того, что водитель автомобиля выезжает на траекторию движения мотоцикла.} \par \textit {Мотоцикл узкий, и водители их не замечают. Большинство аварий происходит, когда водитель не заметил мотоцикл и повернул или выехал перед ним.} \\
\end{minipage}

\begin{minipage}{\linewidth}
\noindent\textbf{81.}  -- \textbf {(D).} Вы собираетесь перестроиться в левую полосу, как показано на рисунке выше. Чтобы заметить автомобили, проезжающие рядом с вами, лучше всего:  \textbf {Повернуть голову влево.} \par \textit {Зеркала не показывают слепую зону, где может скрываться другой автомобиль. Проверяйте слепые зоны поворотом головы (head-checks), прежде чем перестраиваться.} \\
\end{minipage}

\begin{minipage}{\linewidth}
\noindent\textbf{82.}  -- \textbf {(B).} Чтобы подстроиться под дорожное движение и избежать мусора на повороте, показанном на рисунке, следует ехать:  \textbf {В позиции 2.} \par \textit {Центральная часть держит вас подальше от встречного транспорта слева и от песка и мусора возле края дороги справа.} \\
\end{minipage}

\begin{minipage}{\linewidth}
\noindent\textbf{83.}  -- \textbf {(D).} Чтобы пересечь рельсы, идущие параллельно вашему курсу:  \textbf {Отъехать в сторону, затем повернуть и пересечь их под углом не менее 45 градусов.} \par \textit {Пересечение путей опасно. Постепенный переезд под острым углом может привести к попаданию колеса в канавку (паз) и сбросить вас с мотоцикла. Пересекайте их под углом от 45 до 90 градусов. Не выезжайте на встречную полосу или обочину.} \\
\end{minipage}

\begin{minipage}{\linewidth}
\noindent\textbf{84.}  -- \textbf {(A).} На этом рисунке мотоциклист, находящийся в наиболее опасном положении, это:  \textbf {Мотоциклист B.} \par \textit {Мотоциклист B находится рядом с автомобилем на полосе обгона, что является самым уязвимым местом. B находится в слепой зоне автомобиля, и у него меньше всего места для маневра.} \\
\end{minipage}

